\documentclass[conference]{IEEEtran}
\IEEEoverridecommandlockouts
% The preceding line is only needed to identify funding in the first footnote. If that is unneeded, please comment it out.
%Template version as of 6/27/2024

\usepackage{cite}
\usepackage{amsmath,amssymb,amsfonts}
\usepackage{algorithmic}
\usepackage{graphicx}
\usepackage{textcomp}
\usepackage{xcolor}
\usepackage{comment}
\usepackage{booktabs}
\usepackage{tabularx}
\usepackage{url}
\def\BibTeX{{\rm B\kern-.05em{\sc i\kern-.025em b}\kern-.08em
    T\kern-.1667em\lower.7ex\hbox{E}\kern-.125emX}}
\begin{document}

\title{Adapting Audio-JEPA for Low-Resource Thai Automatic Speech Recognition}

\author{\IEEEauthorblockN{Sawit Koseeyaumporn, Kittipong Tapyou, Warapol Pratumta}
\IEEEauthorblockA{\textit{Computer Engineering} \\
\textit{King Mongkut's University of Technology Thonburi}\\
Bangkok, Thailand}
}

\maketitle

\begin{abstract}
Automatic speech recognition for Thai is limited by the scarcity of transcribed speech. Self-supervised pretraining reduces this dependence, but evidence for low-resource recognition comes mainly from contrastive and discrete-unit models. Audio-JEPA learns audio representations by predicting the latent features of masked spectrogram regions, but it was developed for general audio classification, and its coarse patches give too few time steps for character-level transcription. We adapt Audio-JEPA to speech by replacing its square spectrogram patches with a full-frequency projection stem that produces one token every 20 milliseconds, and by masking short temporal spans instead of two-dimensional blocks. The adapted encoder is pretrained on unlabeled Thai speech from the Porjai corpus, either from random initialization or from the original Audio-JEPA weights, and is then fine-tuned with a connectionist temporal classification head. We design a controlled comparison that separates the effects of upstream Audio-JEPA weights, in-domain Thai pretraining, and supervised training alone, using labeled budgets of 1, 10, and 50 hours and the full training set, evaluated with character and word error rates.
\end{abstract}

\section{Introduction}
\label{sec:intro}

Automatic speech recognition (ASR) converts speech into text, but supervised systems typically depend on large amounts of transcribed audio. Accurate transcripts require time and human effort, so this requirement is difficult to meet when labeled speech is scarce \cite{conneau2021xlsr}. Thai is such a case. Paired speech--text data remains limited compared with high-resource languages \cite{shao2025xlsrthai}, and the language brings additional difficulties: it is tonal, distinguishes short and long vowels, and is written without spaces between words. In this work, ``low-resource'' refers primarily to a limited transcription budget. The availability of unlabeled audio and of computation is reported separately.

Conventional ASR describes the speech signal with fixed acoustic features such as mel-frequency cepstral coefficients (MFCCs) or Mel-spectrograms \cite{davis1980mfcc}. Self-supervised learning (SSL) instead learns representations directly from unlabeled audio, which is far cheaper to collect than transcripts. Models such as wav2vec~2.0 and HuBERT pretrain an encoder on untranscribed speech and then adapt it to recognition with few labeled examples \cite{baevski2020wav2vec2,hsu2021hubert}. The question therefore shifts from which features to compute to how an encoder should learn its representations from unlabeled audio.

This paper follows a representation-based approach. Instead of reconstructing the input signal or predicting discrete pseudo-labels, the encoder learns by predicting the \emph{representations} of hidden parts of the input from the visible context. Joint-embedding predictive architectures (JEPAs) implement this principle \cite{lecun2022path,assran2023ijepa}, and Audio-JEPA applies it to audio by predicting latent features of masked Mel-spectrogram patches \cite{tuncay2025audiojepa}. Audio-JEPA learns competitive general-purpose audio representations with a modest training budget. However, it was pretrained on general audio rather than speech, evaluated mainly on classification tasks, and never applied to continuous transcription or to Thai. Its coarse patch grid also produces one time step every 625~ms, which is too coarse for character-level CTC alignment. Whether JEPA-style representation prediction can reduce the labeled-data requirement of Thai ASR is therefore an open question.

We investigate this question with a two-stage framework. First, a speech-adapted Audio-JEPA is pretrained, or further pretrained, on unlabeled Thai speech. Second, its encoder is transferred to a recognizer with a connectionist temporal classification (CTC) head \cite{graves2006ctc} and fine-tuned under controlled labeled-data budgets. The contributions of this work are as follows:
\begin{itemize}
\item We adapt Audio-JEPA to speech by replacing its $16\times16$ spectrogram patches with a full-frequency projection stem that produces one token every 20~ms, and by replacing two-dimensional block masking with short temporal-span masking. All Transformer blocks of the original model can still be transferred.
\item We define a JEPA pretraining and CTC fine-tuning pipeline for Thai, including valid-length handling, CTC feasibility checks, and representation-collapse monitoring.
\item We design a controlled comparison that separates the contribution of upstream Audio-JEPA weights, in-domain Thai JEPA pretraining, and supervised training alone, evaluated with 1, 10, and 50 hours and the full set of labeled Thai speech using character and word error rates.
\end{itemize}

The rest of the paper is organized as follows. Section~\ref{sec:related} reviews related work. Section~\ref{sec:method} describes the proposed method and experimental design. Section~\ref{sec:results} will report the experimental results, and Section~\ref{sec:conclusion} will conclude the paper.

\section{Related Work}
\label{sec:related}

\subsection{Speech Representation Learning}
Early ASR systems represented short speech frames with hand-designed features such as mel-frequency cepstral coefficients or perceptual linear prediction coefficients~\cite{davis1980mfcc,hermansky1990plp}. A hidden Markov model described how speech sounds changed over time, Gaussian mixtures scored the acoustic observations, and a pronunciation dictionary linked sound sequences to words. This separation allowed a vocabulary update without retraining the acoustic model. It also made recognition dependent on the dictionary's coverage and the expertise needed to maintain it~\cite{gales2008hmm}.

Hybrid systems retained the dictionary and decoder while replacing Gaussian mixtures with neural networks. This change reduced word error rates across several recognition tasks, although training typically still required frame-level targets obtained by aligning transcripts to the audio~\cite{hinton2012dnn}. Convolutional networks captured local patterns across frequency bands, while recurrent networks modeled how sounds evolved over time~\cite{abdelhamid2014cnn,graves2013rnn}.

End-to-end recognition simplified supervision by learning from recordings paired with whole transcripts. Connectionist temporal classification (CTC) sums over possible alignments between the speech frames and the transcript~\cite{graves2006ctc}. In Listen, Attend and Spell, an attention-based decoder predicts each character from the recording and the preceding characters, allowing it to spell words outside a pronunciation dictionary~\cite{chan2016las}. The Conformer further improves acoustic modeling by combining local convolutions with attention over longer stretches of speech~\cite{gulati2020conformer}. These approaches remove the need to label individual frames, but their supervised objectives still require a transcript for each training recording.

\subsection{Self-Supervised Learning Paradigms}
Self-supervised learning (SSL) uses information in the recording itself as a training target. Mockingjay, DeCoAR, and TERA recover missing or altered acoustic features, an objective that rewards recovery of background detail as well as speech~\cite{liu2020mockingjay,ling2020decoar,liu2021tera}. Contrastive predictive coding and wav2vec~2.0 learn to distinguish a target representation from sampled alternatives, making the choice of alternatives part of training~\cite{oord2018cpc,baevski2020wav2vec2}. HuBERT and WavLM predict cluster labels at masked positions, which requires a separate clustering step. WavLM also trains on noisy and overlapping speech~\cite{hsu2021hubert,chen2022wavlm}. Data2vec predicts continuous features from a teacher network that has access to the full input~\cite{baevski2022data2vec}.

Joint-embedding predictive architectures also learn from continuous targets. I-JEPA predicts features of a hidden image region from its visible surroundings~\cite{assran2023ijepa}. Applications to spectrograms showed that the choice of which regions to hide affects the resulting audio representations~\cite{fei2023ajepa,riou2024audiojepa}. Tuncay et al.\ pretrained Audio-JEPA on about 5{,}300 hours of AudioSet by predicting a teacher network's features, avoiding both clustering and negative sampling~\cite{tuncay2025audiojepa}. Its evaluation included speaker identification and environmental sound classification. Transcription places a different demand on the encoder. The original patch grid provides only 16 time positions for a 10-second clip, so a character-level CTC decoder cannot align a transcript longer than 16 characters to that grid~\cite{tuncay2025audiojepa,graves2006ctc}.

\subsection{Large-Scale ASR and Speech Language Models}
OpenAI's original Whisper learned from 680{,}000 hours of web audio paired with transcripts or English translations of varying quality. Under this weak supervision, an audio encoder and a text decoder learn jointly to transcribe or translate according to task and language tokens. Whisper could transcribe speech from new datasets without dataset-specific fine-tuning, although recognition error was generally higher for languages with fewer hours of paired audio~\cite{radford2023whisper}.

Speech large language models (LLMs) extend this approach to tasks specified by a written instruction. In Typhoon-Audio, an adapter maps audio encoder outputs into the input space of a pretrained text LLM, much as a vision--language model connects an image encoder to its language model. Changing the prompt lets one model transcribe a spoken question or answer it. This flexibility coexists with substantial differences in transcription accuracy. On Thai Common Voice~17, Manakul et al.\ reported word error rates (WERs) of 95.12\% for Qwen-Audio and 14.17\% for their Thai-adapted Typhoon-Audio~\cite{manakul2025typhoonaudio}.

Recent work at Apple and NVIDIA illustrates how decoding cost and latency shape the choice of recognizer. Hori et al.\ used a text LLM to score candidate transcripts from Apple's experimental recognizer. Delaying these calls improved decoding speed and accuracy on LibriHeavy relative to scoring candidates at every step~\cite{hori2025fusion}. At NVIDIA, Andrusenko et al.\ trained a single transducer for both offline and streaming recognition. Encouraging its predictions from short chunks to agree with those from the full recording improved accuracy when little future audio was available~\cite{andrusenko2026unified}.

\subsection{SSL for Low-Resource and Thai ASR}
SSL can reduce the transcription budget even when pretraining requires substantial audio. With ten minutes of labeled English speech, wav2vec~2.0 achieved 4.8\% WER on LibriSpeech test-clean and 8.2\% on test-other, using about 53{,}000 hours of unlabeled pretraining audio and an external language model~\cite{baevski2020wav2vec2}. XLSR and XLS-R shared a pretrained encoder across languages, and the Massively Multilingual Speech project trained one recognizer for 1{,}107 languages. The shared encoder makes recordings in other languages useful during pretraining, while supervised adaptation supplies examples of how the target language is written~\cite{conneau2021xlsr,babu2022xlsr,pratap2023mms}.

Thai ASR results vary with the type of speech being recognized. On Thai broadcast news, an 18{,}000-word recognizer achieved about 20\% WER on clean, planned speech and just under 30\% across recording conditions~\cite{jongtaveesataporn2008news}. Aung et al.\ found a similar sensitivity to domain after adapting Whisper to Thai. On the Thai Common Voice~13 test set, fine-tuning reduced WER from 12.8\% to 5.5\% for the large model. The same Thonburian Whisper model had 18.7\% WER on financial speech containing English terms~\cite{aung2024thonburian}. Yang et al.\ trained a Thai Zipformer transducer on web audio with automatically generated and refined transcripts. Under the same scoring procedure, both Zipformer and Whisper had higher character error rates (CERs) on GigaSpeech~2 YouTube audio than on Common Voice read speech~\cite{yang2025gigaspeech}.

Reported error rates also depend on decoding and scoring conditions. For streaming recognition, Sirichotedumrong et al.\ reported 14.1\% CER on the Thai online-video benchmark TVSpeech with a 0.6-billion-parameter model and 1.04 seconds of look-ahead~\cite{sirichotedumrong2026streaming}. Even with the model output held fixed, scoring conventions can affect the result. On Thai FLEURS, the CER of Typhoon Whisper Large-v3 fell from 9.98\% to 5.69\% when the reference text followed the team's spelling conventions, including writing out numbers~\cite{sirichotedumrong2026typhoon}. Because Thai spaces do not reliably mark word boundaries, comparing WER across systems also requires consistent word segmentation~\cite{jongtaveesataporn2008news}.

Shao et al.\ provide direct evidence for Thai-specific pretraining. After continuing XLSR pretraining on 36{,}000 hours of Thai audio, including 20{,}000 hours of in-house recordings, they obtained 13.91\% CER on the GigaSpeech~2 test set, compared with 16.74\% for the original encoder under the same CTC fine-tuning procedure~\cite{shao2025xlsrthai}. Speech-specific JEPA studies have also reported ASR results using soft clustering targets~\cite{ioannides2026gmmjepa,ioannides2026sjepa}. An open question is whether predicting continuous features of Thai speech improves transcription when only a few hours have labels. We examine this question by pretraining a speech-adapted Audio-JEPA on Porjai Central Thai recordings~\cite{porjai} and comparing CER under a common CTC decoder at 1, 10, and 50 labeled hours and on the full training set.

\section{Methodology}
\label{sec:method}

\subsection{Research Design}
We propose a two-stage recognizer, denoted TAJA-CTC: self-supervised representation learning from Thai speech, followed by supervised character prediction with connectionist temporal classification (CTC) \cite{graves2006ctc}. Figure~\ref{fig:taja} distinguishes the pretraining graph from the transcription path. This section specifies a proposed implementation; no training outcomes are claimed. The central hypothesis is that masked latent prediction improves label efficiency relative to an identical encoder trained only with CTC.

Audio-JEPA uses an EMA target network and a predictor to learn from masked spectrogram patches \cite{tuncay2025audiojepa}. Its reported frontend has 256 time bins per 10-second clip and a patch width of 16 bins. Thus, we estimate only 16 temporal positions per clip, approximately 625~ms apart. Pooling across frequency would retain this coarse timing, while flattening frequency positions would not create additional acoustic time steps. We therefore replace the two-dimensional patch grid with full-frequency, short-duration speech tokens. This is an adaptation of the Audio-JEPA objective and Transformer backbone, rather than a reproduction of its original frontend.

\begin{figure*}[t]
\centering
\begin{tikzpicture}[
 x=1cm,y=1cm, font=\small,
 box/.style={draw=black!65,rounded corners=2pt,align=center,minimum height=0.82cm,text width=2.50cm,inner sep=5pt},
 online/.style={box,fill=blue!7},
 teacher/.style={box,fill=orange!12},
 loss/.style={box,fill=green!9},
 flow/.style={-{Stealth[length=2mm]},thick},
 note/.style={font=\footnotesize,align=center}
]
\node[anchor=west,font=\bfseries] at (-1.45,0.75) {(a) Self-supervised Thai speech pretraining};
\node[box] (audio) at (0,-0.25) {Unlabeled Thai audio\\16~kHz mono};
\node[box] (mel) at (3.35,-0.25) {80-band log-Mel\\25~ms window / 10~ms hop};
\node[online] (stem) at (6.7,-0.25) {Context projection stem\\2 frames $\rightarrow$ 1 token};
\node[online] (enc) at (10.05,-0.25) {Context Transformer\\visible tokens only};
\node[online] (pred) at (13.4,-0.25) {Predictor + mask tokens\\predict hidden time steps};
\node[teacher] (tstem) at (6.7,-2.4) {Target projection stem\\complete token sequence};
\node[teacher] (tenc) at (10.05,-2.4) {EMA target Transformer\\stop-gradient};
\node[loss] (obj) at (13.4,-2.4) {Normalized latent loss\\masked positions only};
\draw[flow] (audio)--(mel);
\draw[flow] (mel)--(stem);
\draw[flow] (stem)--(enc);
\node[note] at (8.375,0.6) {Mask temporal spans};
\draw[flow] (enc)--(pred);
\draw[flow] (mel.south) |- (tstem.west);
\draw[flow] (tstem)--(tenc);
\draw[flow] (tenc)--(obj);
\draw[flow] (pred)--(obj);
\draw[flow,dashed] (stem)--node[right,note]{EMA weights}(tstem);
\draw[flow,dashed] (enc)--node[right,note]{EMA weights}(tenc);
\node[note,anchor=west] at (-1.45,-3.3) {Gradients update the context stem, context encoder, and predictor; the target branch is updated only by EMA.};
\draw[black!25] (-1.45,-3.8)--(14.9,-3.8);
\node[anchor=west,font=\bfseries] at (-1.45,-4.35) {(b) Supervised adaptation and offline inference};
\node[box] (input) at (0,-5.25) {Thai audio\\same log-Mel frontend};
\node[online] (asr) at (3.35,-5.25) {Retained context stem\\+ Transformer};
\node[online] (ctc) at (6.7,-5.25) {20~ms features\\dropout + linear head};
\node[box] (decode) at (10.05,-5.25) {CTC decoding\\merge repeats, drop blank};
\node[box] (text) at (13.4,-5.25) {Thai transcript\\normalized characters};
\node[loss] (ctcloss) at (6.7,-6.85) {CTC loss\\labeled training only};
\node[box] (labels) at (3.35,-6.85) {Reference transcript\\code-point labels};
\draw[flow] (input)--(asr);
\draw[flow] (asr)--(ctc);
\draw[flow] (ctc)--(decode);
\draw[flow] (decode)--(text);
\draw[flow] (ctc)--(ctcloss);
\draw[flow] (labels)--(ctcloss);
\node[note,text width=5.6cm] at (12,-6.85) {The predictor, target encoder, and training losses are absent during inference.};
\end{tikzpicture}
\caption{Proposed Audio-JEPA adaptation for Thai ASR. Solid arrows show data flow; dashed arrows show exponential moving-average (EMA) parameter updates. The full-frequency speech stem produces an ordered 20~ms sequence. Pretraining predicts hidden latent features; supervised CTC adaptation maps those features to text. All architecture and training choices shown are proposed, not measured results.}
\label{fig:taja}
\end{figure*}


\subsection{Acoustic Frontend and Temporal Encoder}
Audio will be converted to mono at 16~kHz. We will compute 80 log-Mel bands using a 400-sample Hann window, 160-sample hop, 512-point FFT, and a 0--8~kHz Mel filterbank. Power values will be floored at $10^{-10}$ before taking the natural logarithm. Per-band mean and variance will be estimated on training audio only and reused for validation and testing.

Let $X\in\mathbb{R}^{80\times F}$ denote the normalized features. A two-dimensional convolution with kernel and stride $(80,2)$ and 768 output channels will project each non-overlapping pair of frames into one full-frequency token. Equivalently, a learned affine map projects each 160-dimensional frame pair to 768 dimensions. An odd final frame will be zero-padded, giving $T=\lceil F/2\rceil$ tokens at a nominal 20~ms stride; batch padding beyond these valid tokens is excluded by length masks. This explicit frequency-to-feature projection preserves temporal order without averaging the utterance. The 25~ms analysis windows overlap locally; masking therefore removes token inputs, not every waveform sample contributing to neighboring windows.

The context encoder contains 12 pre-layer-normalized Transformer blocks, width 768, 12 attention heads, and feed-forward width 3072. These dimensions permit transfer of compatible Audio-JEPA Transformer blocks. One-dimensional sinusoidal positions replace the original two-dimensional positions. There is no classification token or global temporal pooling. A padding mask excludes invalid time steps from self-attention and all losses. The model is bidirectional and intended for offline transcription; streaming latency is outside this experiment.

\subsection{Thai Speech JEPA Pretraining}
Write $P_\theta$ for the projection stem and $E_\theta$ for the context Transformer. A target copy $(P_{\bar\theta},E_{\bar\theta})$ starts with identical parameters. For each training crop, a temporal mask $M$ will cover approximately 50\% of valid tokens using non-overlapping spans of five tokens, corresponding to 100~ms. A final span may be shortened to meet the mask count. The context Transformer receives only visible tokens with their original positions. The target network receives the complete, uncorrupted sequence.

The predictor projects context outputs to width 384, restores masked positions with learned mask tokens, adds their temporal positions, and processes the ordered sequence using six Transformer blocks with 12 heads and feed-forward width 1536. A final linear projection returns 768-dimensional predictions. Let $z_t$ be the target encoder's last-layer output and $\hat z_t$ the predictor output. Our primary objective is normalized latent regression:
\begin{equation}
\mathcal{L}_{\mathrm{J}}=
\frac{1}{|M|d}\sum_{t\in M}
\left\|\operatorname{LN}(\hat z_t)-
\operatorname{sg}\!\left[\operatorname{LN}(z_t)\right]\right\|_2^2,
\label{eq:jepa}
\end{equation}
where $d=768$, $\operatorname{sg}$ stops gradients, and $\operatorname{LN}$ standardizes each token across channels without learned affine parameters. Normalization and temporal-span masking are proposed speech adaptations; random-token masking is evaluated separately. Only the context branch and predictor receive gradients. After each optimizer update, both the target stem and target Transformer are updated by
\begin{equation}
\bar\theta\leftarrow\tau_s\bar\theta+(1-\tau_s)\theta,
\qquad
\tau_s=1-(1-0.996)\frac{1+\cos(\pi s/S)}{2},
\end{equation}
where $s$ is the update index and $S$ is the pretraining budget. Target dropout is disabled. Padded tokens are excluded from mask sampling and target statistics.

EMA and normalization do not by themselves establish non-collapse. We will record per-channel standard deviations across unpadded target tokens, effective covariance rank, gradient norms, and held-out JEPA loss. Low regression loss accompanied by near-constant embeddings will be reported as a failed representation-learning run. The predeclared acoustic-grounding ablation below tests an additional stabilization mechanism; it is not silently substituted for the primary objective.

\subsection{Initialization and Checkpoint Transfer}
The primary experiment initializes the proposed architecture randomly and pretrains it on the fixed Thai audio pool. A separate transfer experiment will load compatible attention, feed-forward, and layer-normalization weights from an official Audio-JEPA checkpoint, if obtainable from the authors' implementation \cite{audiojepacode}. The new input stem and temporal positions are not loaded from the original patch embedding. The target is copied from the initialized context network, and the predictor is initialized afresh. A parameter-name and shape audit will record every loaded and skipped tensor. Because the public repository does not currently provide a usable checkpoint link, this transfer comparison is conditional; the from-scratch Thai JEPA experiment does not depend on its availability.

\subsection{Supervised CTC Adaptation and Decoding}
After pretraining, the context stem and encoder will be retained; the target branch and predictor will be removed. Unlike the original Audio-JEPA frozen-target evaluation, our downstream choice is the context encoder. A dropout layer with probability 0.1 and a linear classifier produce one distribution over $\mathcal{V}\cup\{\varnothing\}$ per time step, where $\varnothing$ is the CTC blank. For transcript $y$, training minimizes
\begin{equation}
\mathcal{L}_{\mathrm{CTC}}=-\log
\sum_{\pi\in\mathcal{B}^{-1}(y)}
\prod_{t=1}^{T}p(\pi_t\mid x).
\end{equation}
The CTC map $\mathcal{B}$ first merges consecutive identical path labels and then removes blanks \cite{graves2006ctc}. For a target of length $U$ with $R$ adjacent repeated label pairs, we require $T\ge U+R$. Invalid examples will be logged and counted; they will not be concealed by replacing infinite losses with zero. Any training exclusions will be applied to all comparison systems.

Output units are Unicode code points, not word tokens or grapheme clusters. A fixed inventory will include assigned Thai letters, vowels, tone marks and orthographic symbols, ASCII letters and digits, Thai digits, one space token, and an unknown token. Its exact code-point list will be saved before inspecting held-out transcripts. Unsupported symbols map to the unknown token during training; their frequency is reported, and evaluation references retain their normalized characters so that unsupported characters still incur errors.

Text will use Unicode NFC, lowercase Latin letters, and collapsed whitespace. Thai combining vowels and tone marks, numerals, abbreviations, and the repetition mark will be preserved. A versioned list will remove sentence punctuation; no context-dependent number verbalization or repetition expansion is assumed. The same rules apply to labels and decoded text. Primary decoding is greedy CTC. A secondary prefix beam search with beam width 50 and no language model will isolate search effects. No external text or pretrained language model enters the principal results.

\subsection{Optional Acoustic-Grounding Ablation}
Motivated by speech JEPA work using soft clustering anchors \cite{ioannides2026gmmjepa}, we will fit a 100-component diagonal-covariance Gaussian mixture model (GMM) to at most one million normalized training log-Mel frames, sampled with a fixed seed. Adjacent-frame posterior averages provide $q_t$ on the 20~ms grid. An auxiliary linear head on the predictor outputs $r_t$, and training uses
\begin{equation}
\mathcal{L}=\mathcal{L}_{\mathrm{J}}+
\lambda_s\frac{1}{|M|}\sum_{t\in M}
D_{\mathrm{KL}}(q_t\|r_t),
\end{equation}
with $\lambda_s=0.1\max(0,1-2s/S)$. The GMM uses no transcripts or held-out audio and stays fixed. This is our specified auxiliary variant, not an exact reproduction of GMM-JEPA or S-JEPA. Its head is discarded before CTC adaptation.


\section{Experimental Results and Discussion}
\label{sec:results}
% TODO: To be completed after the experiments in Section~\ref{sec:configs} have been run.
% Planned content:
%   - CTC feasibility statistics and data statistics after filtering
%   - Stage 1 pretraining curves and collapse diagnostics
%   - CER/WER for systems A--D and Ref at 1, 10, 50, and ~700 h (mean +/- std over 3 seeds)
%   - Label-efficiency learning curves and error analysis (tone marks, digits, code-switching)
\textit{This section will be completed after the experiments have been conducted.}

\section{Conclusion and Future Work}
\label{sec:conclusion}
% TODO: To be written after the results are available.
\textit{This section will be completed after the experiments have been conducted.}

\section*{Acknowledgment}

The preferred spelling of the word ``acknowledgment'' in America is without 
an ``e'' after the ``g''. Avoid the stilted expression ``one of us (R. B. 
G.) thanks $\ldots$''. Instead, try ``R. B. G. thanks$\ldots$''. Put sponsor 
acknowledgments in the unnumbered footnote on the first page.

\begin{comment}
\section*{References}

Please number citations consecutively within brackets \cite{b1}. The 
sentence punctuation follows the bracket \cite{b2}. Refer simply to the reference 
number, as in \cite{b3}---do not use ``Ref. \cite{b3}'' or ``reference \cite{b3}'' except at 
the beginning of a sentence: ``Reference \cite{b3} was the first $\ldots$''

Number footnotes separately in superscripts. Place the actual footnote at 
the bottom of the column in which it was cited. Do not put footnotes in the 
abstract or reference list. Use letters for table footnotes.

Unless there are six authors or more give all authors' names; do not use 
``et al.''. Papers that have not been published, even if they have been 
submitted for publication, should be cited as ``unpublished'' \cite{b4}. Papers 
that have been accepted for publication should be cited as ``in press'' \cite{b5}. 
Capitalize only the first word in a paper title, except for proper nouns and 
element symbols.

For papers published in translation journals, please give the English 
citation first, followed by the original foreign-language citation \cite{b6}.

\begin{thebibliography}{00}
\bibitem{b1} G. Eason, B. Noble, and I. N. Sneddon, ``On certain integrals of Lipschitz-Hankel type involving products of Bessel functions,'' Phil. Trans. Roy. Soc. London, vol. A247, pp. 529--551, April 1955.
\bibitem{b2} J. Clerk Maxwell, A Treatise on Electricity and Magnetism, 3rd ed., vol. 2. Oxford: Clarendon, 1892, pp.68--73.
\bibitem{b3} I. S. Jacobs and C. P. Bean, ``Fine particles, thin films and exchange anisotropy,'' in Magnetism, vol. III, G. T. Rado and H. Suhl, Eds. New York: Academic, 1963, pp. 271--350.
\bibitem{b4} K. Elissa, ``Title of paper if known,'' unpublished.
\bibitem{b5} R. Nicole, ``Title of paper with only first word capitalized,'' J. Name Stand. Abbrev., in press.
\bibitem{b6} Y. Yorozu, M. Hirano, K. Oka, and Y. Tagawa, ``Electron spectroscopy studies on magneto-optical media and plastic substrate interface,'' IEEE Transl. J. Magn. Japan, vol. 2, pp. 740--741, August 1987 [Digests 9th Annual Conf. Magnetics Japan, p. 301, 1982].
\bibitem{b7} M. Young, The Technical Writer's Handbook. Mill Valley, CA: University Science, 1989.
\bibitem{b8} D. P. Kingma and M. Welling, ``Auto-encoding variational Bayes,'' 2013, arXiv:1312.6114. [Online]. Available: https://arxiv.org/abs/1312.6114
\bibitem{b9} S. Liu, ``Wi-Fi Energy Detection Testbed (12MTC),'' 2023, gitHub repository. [Online]. Available: https://github.com/liustone99/Wi-Fi-Energy-Detection-Testbed-12MTC
\bibitem{b10} ``Treatment episode data set: discharges (TEDS-D): concatenated, 2006 to 2009.'' U.S. Department of Health and Human Services, Substance Abuse and Mental Health Services Administration, Office of Applied Studies, August, 2013, DOI:10.3886/ICPSR30122.v2
\bibitem{b11} K. Eves and J. Valasek, ``Adaptive control for singularly perturbed systems examples,'' Code Ocean, Aug. 2023. [Online]. Available: https://codeocean.com/capsule/4989235/tree
\end{thebibliography}
\end{comment}

\bibliographystyle{IEEEtran}
\bibliography{references}

\end{document}
